% !TEX root = ../../ctfp-print.tex

\lettrine[lhang=0.17]{在}{范畴论中}，似乎一切事物都彼此关联，且每个概念都能从多个角度观察。以\hyperref[products-and-coproducts]{积}的泛构造为例。既然我们对\hyperref[functors]{函子}和\hyperref[natural-transformations]{自然变换}有了更多了解，能否将其简化甚至推广？让我们尝试分析。

\begin{figure}[H]
  \centering
  \includegraphics[width=0.3\textwidth]{images/productpattern.jpg}
\end{figure}

\noindent
积的构造始于两个对象$a$和$b$的选择——这正是我们要构造其积的对象。但"选择对象"具体意味着什么？能否用更范畴论化的语言重新表述这一行为？两个对象构成了一种模式——非常简单的模式。我们可以将这种模式抽象成一个范畴——虽然是个非常简单的范畴，但仍是合法的范畴。我们将这个范畴称为$\cat{2}$，它恰好包含两个对象$1$和$2$，以及各自的恒等态射。现在我们可以将范畴$\cat{C}$中选择两个对象的行为，重新表述为定义一个从范畴$\cat{2}$到$\cat{C}$的函子$D$。函子将对象映射到对象，因此其像就是两个对象（或者可能是同一个对象，若函子发生坍缩，这也是允许的）。它也映射态射——在此情形下只需将恒等态射映射为恒等态射。

这种方法的优势在于建立在范畴论基础之上，摒弃了诸如"选择对象"这类模糊的原始描述（源自人类祖先的采集-狩猎语言体系）。顺便指出，这种方法也易于推广，因为我们完全可以使用比$\cat{2}$更复杂的范畴来定义模式。

\begin{figure}[H]
  \centering
  \includegraphics[width=0.35\textwidth]{images/two.jpg}
\end{figure}

\noindent
继续分析：积定义的下一步是候选对象$c$的选择。同样地，我们可以将这种选择用从单元素范畴出发的函子来表述。实际上，若我们采用Kan扩张方法论，这正是正确的做法。但由于尚未准备讨论Kan扩张，这里可采用另一种技巧：从相同范畴$\cat{2}$到$\cat{C}$的常函子$\Delta$。通过$\Delta_c$即可完成$\cat{C}$中$c$的选择。请记住，$\Delta_c$将所有对象映射为$c$，所有态射映射为$\idarrow[c]$。

\begin{figure}[H]
  \centering
  \includegraphics[width=0.35\textwidth]{images/twodelta.jpg}
\end{figure}

\noindent
现在我们有两个函子$\Delta_c$和$D$在$\cat{2}$与$\cat{C}$之间作用，自然会考虑它们之间的自然变换。由于$\cat{2}$仅有两个对象，自然变换将包含两个分量。对象$1$在$\cat{2}$中被$\Delta_c$映射为$c$，被$D$映射为$a$。因此$\Delta_c$与$D$之间的自然变换在$1$处的分量是从$c$到$a$的态射，记作$p$。类似地，第二个分量是从$c$到$b$的态射$q$——即对象$2$在$\cat{2}$中经$D$的像。这些恰似我们在积原始定义中使用的两个投影。因此我们不必谈论对象和投影的选择，而可以直接讨论函子和自然变换的选择。在这种简单情形下，自然变换的可交换条件自动满足，因为$\cat{2}$中除恒等态射外无其他态射。

\begin{figure}[H]
  \centering
  \includegraphics[width=0.35\textwidth]{images/productcone.jpg}
\end{figure}

\noindent
将此构造推广到包含非平凡态射的范畴$\cat{I}$时，会对$\Delta_c$与$D$之间的变换施加自然性约束。我们将这种变换称为\emph{锥体}（cone），因为$\Delta$的像形成锥体/金字塔的顶点，其边由自然变换的分量构成。$D$的像则构成锥体的底面。

一般而言，构造锥体始于定义模式的范畴$\cat{I}$（通常为小范畴或有限范畴）。选取从$\cat{I}$到$\cat{C}$的函子$D$，称其（或其像）为\emph{图式}（diagram）。选择某个$c$作为锥体顶点，由此定义常函子$\Delta_c$。从$\Delta_c$到$D$的自然变换即为所求锥体。当$\cat{I}$有限时，这对应于连接$c$与图式各节点的一组态射。

\begin{figure}[H]
  \centering
  \includegraphics[width=0.35\textwidth]{images/cone.jpg}
\end{figure}

\noindent
自然性要求该图式中所有三角形（金字塔的侧面）必须可交换。取$\cat{I}$中任意态射$f$，函子$D$将其映射为$\cat{C}$中的态射$Df$，构成某三角形的底边。常函子$\Delta_c$将$f$映射为$c$上的恒等态射，$\Delta$将态射两端压缩为单一对象，此时自然平方图退化为可交换三角形，该三角形的两边即自然变换的分量。

\begin{figure}[H]
  \centering
  \includegraphics[width=0.35\textwidth]{images/conenaturality.jpg}
\end{figure}

\noindent
这就是一个锥体。我们真正关注的是\newterm{泛锥体}（universal cone）——正如在积的定义中选择泛对象那样。

实现途径有多种。例如，我们可以基于给定函子$D$定义\emph{锥体范畴}。该范畴的对象是锥体，但并非$\cat{C}$中每个对象$c$都能成为锥体顶点，因为可能存在从$\Delta_c$到$D$的自然变换。

要构造成范畴还需定义锥体间的态射。这些态射完全由顶点间的态射决定，但并非任意态射都可行。回顾积构造时，我们要求候选对象（顶点）间的态射必须是投影的公共因子。例如：

\src{snippet01}

\begin{figure}[H]
  \centering
  \includegraphics[width=0.35\textwidth]{images/productranking.jpg}
\end{figure}

\noindent
在一般情形下，该条件转化为：包含因子化态射的所有三角形必须可交换。

\begin{figure}[H]
  \centering
  \includegraphics[width=0.35\textwidth]{images/conecommutativity.jpg}
  \caption{连接两个锥体的可交换三角形，其中因子化态射$h$（此处下方锥体是泛锥体，其顶点为$\Lim[D]$）}
\end{figure}

\noindent
我们将这些因子化态射作为锥体范畴中的态射。容易验证这些态射具有合成性，且恒等态射也是因子化态射。因此锥体构成一个范畴。

现在我们可以将泛锥体定义为锥体范畴中的\emph{终对象}。根据终对象定义，存在从任一其他对象到该对象的唯一态射。在此情形下，意味着存在从任意其他锥体顶点到泛锥体顶点的唯一因子化态射。我们将这个泛锥体称为图式$D$的\emph{极限}（limit），记作$\Lim[D]$（文献中常见在$\Lim$符号下方添加指向$I$的左箭头）。简写时常将该锥体的顶点称为极限（或极限对象）。

直观上，极限将整个图式的性质凝聚于单一对象之中。例如，双对象图式的极限就是这两个对象的积。积（连同两个投影）包含了这两个对象的全部信息，而泛性保证了它不含冗余成分。

\section{极限作为自然同构}

关于极限的这种定义仍存在某种不足之处。虽然可行，但我们仍需处理连接任意两个锥体的三角形的可交换条件。若能将其替换为某种自然性条件岂不更为优雅？该如何实现？

我们现在面对的不再是一个锥体，而是整个锥体集合（事实上是一个范畴）。如果极限存在（需要明确说明：这并无保证），其中某个锥体就是泛锥体。对于每个其他锥体，存在唯一因子化态射将其顶点（记作$c$）映射到泛锥体顶点$\Lim[D]$。（实际上，"其他"这个词可以省略，因为恒等态射将泛锥体映射到自身且显然可因子化。）重复强调重点：给定任意锥体，存在唯一特殊类型的态射。我们建立了从锥体到特殊态射的一一对应。

这个特殊态射属于同态集$\cat{C}(c, \Lim[D])$。该同态集的其他成员不具备因子化两个锥体映射的特性。我们的目标是对每个$c$，从$\cat{C}(c, \Lim[D])$中选择一个满足特定可交换条件的态射。这听起来是否像是在定义自然变换？确实如此！

但究竟是哪两个函子通过这种变换相关联？

第一个函子将$c$映射到集合$\cat{C}(c, \Lim[D])$。这是从$\cat{C}$到$\Set$的函子——将对象映射到集合。事实上这是一个逆变函子。其在态射上的作用定义如下：取从$c'$到$c$的态射：
$$f \Colon c' \to c$$
该函子将$c'$映射到集合$\cat{C}(c', \Lim[D])$。要定义该函子在$f$上的作用（即提升$f$），我们需要定义$\cat{C}(c, \Lim[D])$与$\cat{C}(c', \Lim[D])$之间的对应映射。取$\cat{C}(c, \Lim[D])$的一个元素$u$：
$$u \Colon c \to \Lim[D]$$
将$u$与$f$前复合得到：
$$u \circ f \Colon c' \to \Lim[D]$$
这正是$\cat{C}(c', \Lim[D])$的元素——因此我们找到了态射的对应关系：

\src{snippet02}
注意逆变函子特有的$c$与$c'$顺序反转。

\begin{figure}[H]
  \centering
  \includegraphics[width=0.35\textwidth]{images/homsetmapping.jpg}
\end{figure}

\noindent
要定义自然变换，我们需要另一个从$\cat{C}$到$\Set$的函子。这次我们将考虑锥体的集合。锥体本质上是自然变换，因此我们考察自然变换集$\mathit{Nat}(\Delta_c, D)$。从$c$到该自然变换集的映射是一个（逆变）函子。如何证明这点？同样定义其在态射上的作用：
$$f \Colon c' \to c$$
提升$f$应给出从$\cat{I}$到$\cat{C}$的两个函子之间的自然变换映射：
$$\mathit{Nat}(\Delta_c, D) \to \mathit{Nat}(\Delta_{c'}, D)$$
如何映射自然变换？每个自然变换都是分量（对应$\cat{I}$的每个元素）的选取。某个$\alpha \in \mathit{Nat}(\Delta_c, D)$在对象$a$处的分量是态射：
$$\alpha_a \Colon \Delta_c a \to D a$$
根据常函子$\Delta$的定义，即：
$$\alpha_a \Colon c \to D a$$
给定$f$和$\alpha$，我们需要构造$\beta \in \mathit{Nat}(\Delta_{c'}, D)$。其在$a$处的分量应为态射：
$$\beta_a \Colon c' \to D a$$
通过将$\alpha_a$与$f$前复合即可获得：
$$\beta_a = \alpha_a \circ f$$
容易证明这些分量确实构成自然变换。

\begin{figure}[H]
  \centering
  \includegraphics[width=0.4\textwidth]{images/natmapping.jpg}
\end{figure}

\noindent
给定态射$f$，我们逐分量构建了两个自然变换之间的映射。该映射定义了函子：
$$c \to \mathit{Nat}(\Delta_c, D)$$
的\code{contramap}。我刚刚展示的是：我们有两个从$\cat{C}$到$\Set$的（逆变）函子。未做任何假设——这些函子始终存在。

顺便提及，第一个函子在范畴论中扮演重要角色，当我们讨论Yoneda引理时会再次遇到。从任意范畴$\cat{C}$到$\Set$的逆变函子被称为"预层"（presheaf）。这个特定函子称为\newterm{可表预层}（representable presheaf）。第二个函子也是预层。

现在有了两个函子，就可以讨论它们之间的自然变换。无需赘述，直接得出结论：函子$D:\cat{I}\to\cat{C}$具有极限$\Lim[D]$当且仅当存在上述两个函子之间的自然同构：
$$\cat{C}(c, \Lim[D]) \simeq \mathit{Nat}(\Delta_c, D)$$
提醒您自然同构的定义：其每个分量都是同构态射（即可逆态射）。

此处不提供该命题的完整证明。过程虽繁琐但直接。处理自然变换时通常聚焦于分量（即态射）。在此情形下，由于两个函子的目标都是$\Set$，自然同构的分量将是函数。这些高阶函数将同态集映射到自然变换集。通过分析函数对其参数的作用（此处参数是$\cat{C}(c, \Lim[D])$的态射，结果是$\mathit{Nat}(\Delta_c, D)$的自然变换，即锥体），可以追踪到底层态射的对应关系从而完成证明。

最重要结论是：该同构的自然性条件正好对应锥体映射的可交换条件。

预告未来的内容：集合$\mathit{Nat}(\Delta_c, D)$可视为函子范畴中的同态集；因此我们的自然同构关联了两个同态集，这指向更一般的伴随关系(adjunction)。

\section{极限的例子}

我们已经看到，范畴积是通过一个简单的范畴$\cat{2}$生成的图示的极限。

还有一个更简单的极限例子：终对象。第一反应可能是认为单对象范畴会导致终对象，但实际情况更为极端：终对象是由空范畴生成的极限。从空范畴出发的函子不会选择任何对象，因此锥体退化为仅剩顶点。泛锥就是那个具有唯一态射从任意其他顶点指向它的特殊顶点。这正是终对象的定义。

下一个有趣的极限称为\emph{等化子（equalizer）}。它由两个对象组成的范畴生成，该范畴包含两个平行态射（以及恒等态射）。这个范畴选择了一个$\cat{C}$中的图示，包含两个对象$a$和$b$，以及两个态射：

\src{snippet03}

要构建这个图示上的锥体，我们需要添加顶点$c$和两个投影：

\src{snippet04}

\begin{figure}[H]
  \centering
  \includegraphics[width=0.35\textwidth]{images/equalizercone.jpg}
\end{figure}

\noindent
这两个三角形必须可交换：

\src{snippet05}

这表明$q$由其中一个方程唯一确定，例如\code{q = f . p}，因此可以从图中省略。于是只剩下一个条件：

\src{snippet06}

直观理解的话，在$\Set$范畴中，函数$p$的像会选择$a$的一个子集。当限制在这个子集上时，函数$f$和$g$相等。

例如，取$a$为由坐标$x$和$y$参数化的二维平面，取$b$为实数直线，并令：

\src{snippet07}

这两个函数的等化子就是实数集合（顶点$c$）和函数：

\src{snippet08}

注意$(p~t)$在二维平面上定义了一条直线。沿这条直线，两个函数相等。

当然还有其他集合$c'$和函数$p'$也可能满足：

\src{snippet09}

但它们都会通过$p$唯一分解。例如，我们可以取单元素集$\cat{()}$作为$c'$和函数：

\src{snippet10}

这是一个良好锥体，因为$f (0, 0) = g (0, 0)$。但它不是泛的，因为存在通过$h$的唯一分解：

\src{snippet11}

其中

\src{snippet12}

\begin{figure}[H]
  \centering
  \includegraphics[width=0.35\textwidth]{images/equilizerlimit.jpg}
\end{figure}

\noindent
因此等化子可用于求解$f~x = g~x$类型的方程。但它的意义更广泛，因为它通过对象和态射而非代数方法定义。

另一个极限---拉回（pullback）体现了更一般的方程求解思想。此时我们仍想"等化"两个态射，但这次它们的定义域不同。我们从形状为$1\rightarrow2\leftarrow3$的三对象范畴开始。对应的图示包含三个对象$a$、$b$、$c$和两个态射：

\src{snippet13}

这个图示常被称为\emph{共跨度（cospan）}。

在此图示上构建的锥体包含顶点$d$和三个态射：

\src{snippet14}

\begin{figure}[H]
  \centering
  \includegraphics[width=0.35\textwidth]{images/pullbackcone.jpg}
\end{figure}

\noindent
可交换性条件告诉我们$r$完全由其他态射确定，可以从图中省略。因此只剩下：

\src{snippet15}
拉回就是这种形状的泛锥体。

\begin{figure}[H]
  \centering
  \includegraphics[width=0.35\textwidth]{images/pullbacklimit.jpg}
\end{figure}

\noindent
再次聚焦到集合范畴，可以认为对象$d$由那些使得$f$作用于第一个分量等于$g$作用于第二个分量的$a$和$c$的元素对组成。如果这仍过于抽象，考虑$g$为常数函数的情况，例如$g~\_ = 1.23$（假设$b$是实数集）。此时实际在求解方程：

\src{snippet16}

在这种情况下，只要$c$非空，其具体选择无关紧要，可以取单元素集。例如$a$可以是三维向量集，$f$为向量长度。此时拉回就是所有形如$(v, ())$的对，其中$v$是长度为1.23的向量（即方程$\sqrt{(x^{2}+y^{2}+z^{2})} = 1.23$的解），$()$是单元素集的哑元素。

但拉回在编程中也有更广泛的应用。例如，将C++类视为以继承关系为态射的范畴。我们视继承为传递关系，若\code{C}继承自\code{B}且\code{B}继承自\code{A}，则\code{C}也继承自\code{A}。同时假设每个类都继承自身，从而保证恒等箭头。这种设计使继承与子类型对齐。C++支持多重继承，可以构造菱形继承图：类\code{B}和\code{C}都继承自\code{A}，而类\code{D}多重继承自\code{B}和\code{C}。通常\code{D}会包含两份\code{A}的副本，但可通过虚继承解决这个问题。

在这个图示中要求\code{D}成为拉回意味着什么？这意味着任何从\code{B}和\code{C}多重继承的类\code{E}都是\code{D}的子类。这在名义子类型（nominal subtyping）的C++中无法直接表达（编译器不会推导这类类关系---需要"鸭子类型"）。但我们可以通过判断从\code{E}到\code{D}的类型转换是否安全来间接验证。若\code{D}只是\code{B}和\code{C}的基础组合（无额外数据和方法覆盖），则此转换安全。当然，若\code{B}和\code{C}的方法存在命名冲突，则不存在这样的拉回。

\begin{figure}[H]
  \centering
  \includegraphics[width=0.25\textwidth]{images/classes.jpg}
\end{figure}

\noindent
拉回在类型推断中还有更高级的应用。常常需要对两个表达式的类型进行\emph{统一（unify）}。例如，编译器要推断函数：

\begin{snip}{haskell}
twice f x = f (f x)
\end{snip}
会为所有变量和子表达式分配初步类型：
\begin{snip}{haskell}
f       :: t0
x       :: t1
f x     :: t2
f (f x) :: t3
\end{snip}
由此推导出：
\begin{snip}{haskell}
twice :: t0 -> t1 -> t3
\end{snip}
同时还会得到函数应用规则产生的约束条件：
\begin{snip}{haskell}
t0 = t1 -> t2 -- 因为 f 被应用于 x
t0 = t2 -> t3 -- 因为 f 被应用于 (f x)
\end{snip}
这些约束需要通过找到一组类型（或类型变量）进行统一，使得替换到两个表达式中后产生相同类型。一种可能的替换是：
\begin{snip}{haskell}
t1 = t2 = t3 = Int
twice :: (Int -> Int) -> Int -> Int
\end{snip}
但这显然不是最通用的。最通用的替换通过拉回获得。这里不展开细节（超出本书范围），但你可以自行验证结果应该是：
\begin{snip}{haskell}
twice :: (t -> t) -> t -> t
\end{snip}
其中\code{t}是一个自由类型变量。

\section{余极限}

正如范畴论中的所有构造一样，极限在对偶范畴中具有其对偶图像。当你将锥体(cone)中所有箭头的方向反转后，会得到一个上锥(co-cone)，其中的泛对象称为余极限(colimit)。需要注意的是这种反转也会影响分解态射(factorizing morphism)，现在它从泛上锥流向任意其他上锥。

\begin{figure}[H]
  \centering
  \includegraphics[width=0.35\textwidth]{images/colimit.jpg}
  \caption{通过分解态射$h$连接两个顶点的上锥。}
\end{figure}

\noindent
余极限的一个典型例子是余积(coproduct)，它对应由范畴$\cat{2}$生成的图示——这个范畴正是我们在定义乘积时所使用的。

\begin{figure}[H]
  \centering
  \includegraphics[width=0.35\textwidth]{images/coproductranking.jpg}
\end{figure}

\noindent
乘积与余积各自以不同方式体现了对象对(object pair)的本质。

正如终对象(terminal object)是一个极限，始对象(initial object)则是对应于空范畴图示的余极限。

拉回(pullback)的对偶概念称为\emph{推出(pushout)}。它基于由范畴$1\leftarrow2\rightarrow3$生成的称为"span"的图示。

\section{连续性}

之前我曾说过，函子接近于范畴上的连续映射的概念，因为它们从不破坏已有的连接（态射）。实际上，从范畴$\cat{C}$到$\cat{C'}$的\emph{连续函子}(continuous functor) $F$的正式定义包含要求其保持极限的条件。每个$\cat{C}$中的图示$D$都可以通过简单组合两个函子映射到$\cat{C'}$中的图示$F \circ D$。$F$的连续性条件表明：若图示$D$具有极限$\Lim[D]$，则图示$F \circ D$也具有极限，且该极限等于$F (\Lim[D])$。

\begin{figure}[H]
  \centering
  \includegraphics[width=0.6\textwidth]{images/continuity.jpg}
\end{figure}

\noindent
注意到由于函子将态射映射到态射，并将复合映射到复合，因此锥体的像总是锥体。一个交换三角形必然被映射为另一个交换三角形（函子保持复合关系）。对于分解态射同样成立：分解态射的像也是分解态射。因此每个函子都\emph{几乎}连续。可能出现问题的是唯一性条件。在$\cat{C'}$中的分解态射可能不唯一，也可能存在其他在$\cat{C}$中不可用的"更优锥体"。

Hom-函子是连续函子的一个典型例子。回忆一下，Hom-函子$\cat{C}(a, b)$在其第一个变量上是逆变的，在第二个变量上是协变的。换句话说，它是如下函子：
\[\cat{C}^\mathit{op} \times \cat{C} \to \Set\]
当固定第二个参数时，Hom-集合函子（此时成为可表预层）会将$\cat{C}$中的余极限映射到$\Set$中的极限；而当固定第一个参数时，则将极限映射到极限。

在Haskell中，Hom-函子表现为任意两种类型到函数类型的映射，即参数化函数类型。当我们固定第二个参数（例如设为\code{String}）时，就得到逆变函子：

\src{snippet17}
连续性意味着当\code{ToString}作用于余极限（例如余积\code{Either b c}）时，会产生极限；在此情形下即两个函数类型的乘积：

\src{snippet18}
确实，任何接受\code{Either b c}作为输入的函数本质上都是通过case语句实现的，其中两种情况分别由一对函数处理。

类似地，当我们固定Hom-集合的第一个参数时，就得到熟悉的Reader函子。其连续性意味着：任一返回乘积类型的函数等价于两个函数的乘积；特别地：

\src{snippet19}
我知道你在想什么：你根本不需要范畴论来理解这些概念。没错！但我仍然惊叹于这些结果能够从第一原理推导出来，完全无需涉及比特与字节、处理器架构、编译器技术甚至lambda演算。

如果你好奇"极限"和"连续性"这些名称的来源，它们源自微积分中相应的概念。在分析学中，极限和连续性是通过开邻域定义的。定义拓扑的开集本身构成一个范畴（偏序集）。

\section{挑战}

\begin{enumerate}
  \tightlist
  \item
        如何描述C++类范畴中的推出(pushout)？
  \item
        证明恒等函子(identity functor)
        $\mathbf{Id} \Colon \cat{C} \to \cat{C}$ 的极限(limit)是初始对象(initial object)。
  \item
        给定集合的子集构成一个范畴(category)。该范畴中的态射(morphism)定义为：
        当第一个集合是第二个集合的子集时，存在连接这两个集合的箭头(arrows)。
        在此类别中两个集合的拉回(pullback)是什么？推出(pushout)是什么？
        初始对象(initial object)和终对象(terminal object)分别是什么？
  \item
        你能推测余等化子(coequalizer)的定义吗？
  \item
        证明在具有终对象(terminal object)的范畴中，向终对象的拉回(pullback)就是积(product)。
  \item
        同理，证明从初始对象(initial object)（若存在）的推出(pushout)就是余积(coproduct)。
\end{enumerate}